\chapter{Glosario de términos}
\label{ap:glosario}

Definiciones breves de los conceptos centrales, para consulta rápida.

\begin{description}[style=nextline,leftmargin=0pt]
  \item[Beneficio ($\Pi$)] Diferencia entre el ingreso total y el costo total: $\Pi = I - C$.
  \item[Ceteris paribus] «Manteniendo todo lo demás constante». Supuesto para aislar el efecto de una sola variable.
  \item[Costo marginal ($\cmg$)] Costo de producir una unidad adicional; derivada del costo total.
  \item[Costo medio ($\cme$)] Costo total dividido por la cantidad: costo por unidad.
  \item[Demanda inversa] Función que da el precio en función de la cantidad, $p(Q)$; despeje de la demanda.
  \item[Elasticidad ($\varepsilon$)] Variación porcentual de una variable ante la variación porcentual de otra; medida adimensional de sensibilidad.
  \item[Equilibrio de mercado] Precio y cantidad donde la cantidad demandada iguala a la ofrecida.
  \item[Equilibrio de Nash] Situación en que ninguna parte gana desviándose, dadas las decisiones de las demás.
  \item[Estática comparativa] Comparación de dos equilibrios antes y después de cambiar un parámetro.
  \item[Excedente del consumidor (EC)] Diferencia entre lo que los consumidores estaban dispuestos a pagar y lo que pagaron.
  \item[Excedente del productor (EP)] Diferencia entre el precio recibido y el costo marginal de cada unidad.
  \item[Ingreso marginal ($\img$)] Ingreso adicional por vender una unidad más; derivada del ingreso total.
  \item[Matriz de Leontief] $(I-A)^{-1}$: requerimientos totales (directos e indirectos) de producción por unidad de demanda final.
  \item[Mediana] Valor que deja la mitad de los datos por debajo y la mitad por encima; medida robusta.
  \item[Optimización] Búsqueda del valor de una variable que maximiza o minimiza una función objetivo.
  \item[Payback (período de recupero)] Tiempo necesario para recuperar la inversión inicial con los flujos del proyecto.
  \item[Programación lineal] Optimización de un objetivo lineal sujeto a restricciones lineales.
  \item[TIR] Tasa de descuento que hace el VAN igual a cero; rentabilidad propia del proyecto.
  \item[Valor tiempo del dinero] Principio de que un peso hoy vale más que un peso en el futuro.
  \item[VAN] Valor Actual Neto: suma de los flujos de un proyecto descontados al presente.
\end{description}
